\chapter{Exponentielle et logarithme}\label{ch:g12:exp}

La \omterm{thm:g12:exp:existence}{fonction exponentielle} est l'unique \omterm{def:g11:func:function}{fonction} égale à sa propre \omterm{def:g12:deriv:derivative}{dérivée}
et valant $1$ en $0$. Elle convertit les sommes en produits~; son inverse,
le \omterm{def:g12:exp:ln}{logarithme népérien}, convertit les produits en sommes. Ensemble, ils
décrivent tout phénomène dont le taux de variation est proportionnel à sa
taille~: désintégration radioactive, croissance de population, intérêts
composés.

\section{La fonction exponentielle}

\begin{theorem}[Existence et unicité]\label{thm:g12:exp:existence}
Il existe une unique \omterm{def:g11:func:function}{fonction} \omterm{def:g12:deriv:derivative}{dérivable} $f \colon \R \to \R$ telle que
\[
f' = f \qquad\text{et}\qquad f(0) = 1 .
\]
On l'appelle la \emph{fonction exponentielle}\index{exponentielle} et on
l'écrit $\exp$, ou $x \mapsto \eu^x$.
\end{theorem}

\begin{proof}[Démonstration de l'unicité]
D'abord, une telle \omterm{def:g11:func:function}{fonction} ne s'annule jamais. En effet, soit
$g(x) = f(x)f(-x)$~; alors
\[
g'(x) = f'(x)f(-x) - f(x)f'(-x) = f(x)f(-x) - f(x)f(-x) = 0,
\]
donc $g$ est constante égale à $g(0) = 1$~: pour tout $x$,
$f(x)f(-x) = 1$, et en particulier $f(x) \neq 0$.

Soient maintenant $f_1, f_2$ deux solutions et $h = \dfrac{f_1}{f_2}$
(légitime car $f_2$ ne s'annule jamais). Alors
\[
h' = \frac{f_1' f_2 - f_1 f_2'}{f_2^2} = \frac{f_1 f_2 - f_1 f_2}{f_2^2} = 0,
\]
donc $h$ est constante égale à $h(0) = 1$, \emph{c'est-à-dire}
$f_1 = f_2$.

L'\emph{existence} est admise à ce niveau (on peut l'obtenir via la \omterm{def:g12:seq:sequence}{suite}
de l'\cref{exo:g12:seq:10}, ou comme inverse du \omterm{def:g12:exp:ln}{logarithme} construit par
intégration dans le \cref{ch:g12:integ}).
\end{proof}

\begin{proposition}[Équation fonctionnelle]\label{prop:g12:exp:funceq}
Pour tous $x, y \in \R$ et $n \in \Z$~:
\[
\eu^{x+y} = \eu^x\,\eu^y, \qquad
\eu^{-x} = \frac{1}{\eu^x}, \qquad
\eu^{x-y} = \frac{\eu^x}{\eu^y}, \qquad
\bigl(\eu^{x}\bigr)^n = \eu^{nx}.
\]
De plus $\eu^x > 0$ pour tout $x$.
\end{proposition}

\begin{proof}
Fixer $y$ et considérer $\varphi(x) = \dfrac{\exp(x+y)}{\exp(x)\exp(y)}$.
Son numérateur et son dénominateur, comme \omterm{def:g11:func:function}{fonctions} de $x$, sont tous deux
des solutions de $f' = f$ à la constante $\exp(y)$ près~; en dérivant
$\varphi$ directement (règle du quotient) on obtient $\varphi' = 0$, donc
$\varphi \equiv \varphi(0) = 1$, prouvant la première identité. Prendre
$y = -x$ donne la deuxième (avec la valeur $\eu^0 = 1$), et la troisième
suit. La quatrième est une récurrence à partir de la première pour
$n \geq 0$, étendue à $n < 0$ par la deuxième.

Positivité~: $\eu^x = \left(\eu^{x/2}\right)^2 \geq 0$ et $\eu^x \neq 0$
(montré dans le \cref{thm:g12:exp:existence}), donc $\eu^x > 0$.
\end{proof}

\begin{proposition}[Variations et limites]\label{prop:g12:exp:variations}
L'exponentielle est strictement \omterm{def:g12:seq:monotonic}{croissante}, \omterm{def:g12:deriv:convex}{convexe}, et
\[
\lim_{x\to-\infty} \eu^x = 0, \qquad
\lim_{x\to+\infty} \eu^x = +\infty .
\]
\end{proposition}

\begin{proof}
$(\exp)' = \exp > 0$ donne la croissance stricte~; $(\exp)'' = \exp > 0$
donne la \omterm{def:g12:deriv:convex}{convexité}. Par \omterm{def:g12:deriv:convex}{convexité}, $\eu^x \geq 1 + x$ (\omterm{def:g12:deriv:tangent}{tangente} en $0$,
voir l'\cref{ex:g12:deriv:ineq}), donc $\eu^x \to +\infty$ lorsque
$x \to +\infty$ par comparaison. Puis $\eu^{x} = 1/\eu^{-x} \to 0$
lorsque $x \to -\infty$.
\end{proof}

\begin{omfigure}
\begin{tikzpicture}
\begin{axis}[omaxis,
  width=8.8cm, height=6cm,
  xmin=-3.3, xmax=2.4, ymin=-2.4, ymax=7.8,
  xtick={-2,-1,1,2}, ytick={1}]
  \addplot[omProp, thick, domain=-3.2:2.3] {1 + x};
  \addplot[omDef, thick, domain=-3.2:2.0] {exp(x)};
  \fill (axis cs:0,1) circle (1.5pt);
  \node[omDef, font=\small, anchor=east] at (axis cs:1.35,5.2) {$\eu^x$};
  \node[omProp, font=\small, anchor=north west] at (axis cs:1.05,1.6) {$y = 1+x$};
\end{axis}
\end{tikzpicture}

{\small L'exponentielle~: strictement \omterm{def:g12:seq:monotonic}{croissante}, \omterm{def:g12:deriv:convex}{convexe}, avec
$\lim_{-\infty} \eu^x = 0$. Elle est au-dessus de sa \omterm{def:g12:deriv:tangent}{tangente} en $0$~:
$\eu^x \geq 1 + x$.}
\end{omfigure}

\begin{theorem}[Comparaison de croissances]\label{thm:g12:exp:growth}
\index{comparaison de croissances}
Pour tout entier $n \geq 1$,
\[
\lim_{x\to+\infty} \frac{\eu^x}{x^n} = +\infty,
\qquad
\lim_{x\to-\infty} x^n \eu^x = 0 .
\]
En d'autres termes~: l'exponentielle bat toute puissance de $x$.
\end{theorem}

\begin{proof}
Pour $n = 1$~: en appliquant $\eu^t \geq 1 + t \geq t$ en $t = x/2$,
\[
\frac{\eu^x}{x} = \frac{\left(\eu^{x/2}\right)^2}{x}
\geq \frac{(x/2)^2}{x} = \frac{x}{4} \xrightarrow[x\to+\infty]{} +\infty .
\]
Pour $n$ général, écrire
\[
\frac{\eu^x}{x^n}
= \frac{1}{n^n}\left(\frac{\eu^{x/n}}{x/n}\right)^{n}
\xrightarrow[x\to+\infty]{} +\infty
\]
par le cas $n=1$ et composition. La limite en $-\infty$ suit par le
changement $x \mapsto -x$~:
$\abs{x^n \eu^x} = \frac{\abs{x}^n}{\eu^{\abs x}} \to 0$.
\end{proof}

\begin{omfigure}
\begin{tikzpicture}
\begin{axis}[omaxis,
  width=8.8cm, height=6cm,
  xmin=-0.4, xmax=5.6, ymin=-12, ymax=155,
  xtick={1,2,3,4,5}, ytick={50,100,150}]
  \addplot[omThm, thick, domain=0:5.4] {x^2};
  \addplot[omProp, thick, domain=0:5.3] {x^3};
  \addplot[omDef, thick, domain=-0.4:5.0] {exp(x)};
  \fill (axis cs:4.536,93.4) circle (1.5pt);
  \draw[black, dashed, thin] (axis cs:4.536,0) -- (axis cs:4.536,93.4);
\end{axis}
\end{tikzpicture}

{\small L'exponentielle (bleu) finit par battre toute puissance~: elle
dépasse $x^2$ (rouge) tôt et $x^3$ (orange) en $x \approx 4.54$.}
\end{omfigure}

\section{Le logarithme népérien}

\begin{definition}[Logarithme népérien]\label{def:g12:exp:ln}
L'exponentielle est \omterm{def:g12:limcont:continuity}{continue} et strictement \omterm{def:g12:seq:monotonic}{croissante} de $\R$ sur
$\intoo{0}{+\infty}$~; par le théorème de la bijection
(\cref{thm:g12:limcont:bijection}), pour tout $y > 0$ l'\omterm{def:g10:algebra:equation}{équation}
$\eu^x = y$ a une unique solution. Cette solution est le \emph{logarithme
népérien}\index{logarithme} de $y$, noté $\ln y$. Ainsi
\[
\text{pour } x \in \R,\ y > 0: \qquad y = \eu^x \iff x = \ln y .
\]
En particulier $\ln 1 = 0$, $\ln \eu = 1$, et
$\eu^{\ln y} = y$, $\ln(\eu^x) = x$.
\end{definition}

\begin{omfigure}
\begin{tikzpicture}
\begin{axis}[omaxis,
  width=8.6cm, axis equal image,
  xmin=-4.3, xmax=6.3, ymin=-4.3, ymax=6.3,
  xtick={1}, ytick={1}]
  \addplot[black, dashed, domain=-3.9:5.9] {x};
  \addplot[omDef, thick, domain=-4.1:1.79] {exp(x)};
  \addplot[omThm, thick, domain=0.0167:5.9, samples=120] {ln(x)};
  \node[omDef, font=\small, anchor=east] at (axis cs:0.95,2.6) {$\eu^x$};
  \node[omThm, font=\small, anchor=north] at (axis cs:4.3,1.25) {$\ln x$};
  \node[font=\small, anchor=west, rotate=45] at (axis cs:4.6,4.85) {$y=x$};
\end{axis}
\end{tikzpicture}

{\small Les courbes de $\exp$ et $\ln$ sont \omterm{def:g10:functions:function}{images} miroir l'une de l'autre
par rapport à la droite $y = x$~: le point $(x, \eu^x)$ se reflète en
$(\eu^x, x)$.}
\end{omfigure}

\begin{proposition}[Propriétés algébriques]\label{prop:g12:exp:lnalg}
Pour tous $a, b > 0$ et $n \in \Z$~:
\[
\ln(ab) = \ln a + \ln b, \quad
\ln\frac{1}{a} = -\ln a, \quad
\ln\frac{a}{b} = \ln a - \ln b, \quad
\ln(a^n) = n \ln a, \quad
\ln\sqrt{a} = \tfrac12 \ln a .
\]
\end{proposition}

\begin{proof}
$\eu^{\ln a + \ln b} = \eu^{\ln a}\eu^{\ln b} = ab$, et en prenant $\ln$
des deux membres on obtient la première identité. Les autres suivent par
le même mécanisme des identités correspondantes de la
\cref{prop:g12:exp:funceq}.
\end{proof}

\begin{proposition}[Propriétés analytiques]\label{prop:g12:exp:lnanalytic}
La \omterm{def:g11:func:function}{fonction} $\ln$ est \omterm{def:g12:deriv:derivative}{dérivable} sur $\intoo{0}{+\infty}$ avec
\[
(\ln x)' = \frac{1}{x},
\]
strictement \omterm{def:g12:seq:monotonic}{croissante}, \omterm{def:g12:deriv:convex}{concave}, et
$\lim\limits_{x\to0^+} \ln x = -\infty$,
$\lim\limits_{x\to+\infty} \ln x = +\infty$. De plus, pour tout entier
$n \geq 1$,
\[
\lim_{x\to+\infty} \frac{\ln x}{x^{1/n}} = 0, \qquad
\lim_{x\to0^+} x \ln x = 0 .
\]
\end{proposition}

\begin{proof}
La dérivabilité de la \omterm{def:g11:func:function}{fonction} inverse est admise à ce niveau~; en l'accordant,
dériver l'identité $\eu^{\ln x} = x$ par la règle de la chaîne~:
$(\ln x)'\,\eu^{\ln x} = 1$, donc $(\ln x)' = \frac{1}{x} > 0$, d'où la
croissance stricte, et $(\ln x)'' = -\frac1{x^2} < 0$, d'où la concavité.
Les limites en $0^+$ et $+\infty$ reflètent celles de $\exp$ via la
bijection. Pour la \omterm{thm:g12:exp:growth}{comparaison de croissances}, substituer $x = \eu^{t}$~:
$\frac{\ln x}{x} = \frac{t}{\eu^t} \to 0$ lorsque $t \to +\infty$ par le
\cref{thm:g12:exp:growth}, et de même pour les autres limites avec
$x = \eu^{-t}$, donnant $x\ln x = -t\eu^{-t} \to 0$.
\end{proof}

\begin{method}[Résoudre des équations avec $\exp$ et $\ln$]\label{met:g12:exp:equations}
Les deux \omterm{def:g11:func:function}{fonctions} sont strictement \omterm{def:g12:seq:monotonic}{croissantes}, donc on peut les appliquer
à (ou les retirer de) les deux membres d'une \omterm{def:g10:algebra:equation}{équation} ou d'une inégalité
sans changer le sens~:
\[
\eu^{u} = \eu^{v} \iff u = v, \qquad
\eu^{u} \leq \eu^{v} \iff u \leq v,
\]
et de même pour $\ln$ sur des quantités positives. \emph{Toujours vérifier
les \omterm{def:g11:func:function}{ensembles de définition} d'abord} ($\ln$ exige des arguments positifs).
Pour des \omterm{def:g10:algebra:equation}{équations} du type $a^x = b$ avec $a>0$, $a \neq 1$, réécrire
$a^x = \eu^{x\ln a}$ et résoudre $x = \frac{\ln b}{\ln a}$.
\end{method}

\begin{example}\label{ex:g12:exp:halflife}
Un \omterm{def:g10:proba:sample}{échantillon} radioactif se désintègre suivant $N(t) = N_0\,\eu^{-\lambda t}$.
Sa \emph{demi-vie} $T$ satisfait $N(T) = N_0/2$, \emph{c'est-à-dire}
$\eu^{-\lambda T} = \frac12$, donc $T = \frac{\ln 2}{\lambda}$~: la demi-vie
ne dépend pas de la quantité initiale.
\end{example}

\section{Exercices}

\begin{exercise}[$\star$]\label{exo:g12:exp:1}
Simplifier
$\dfrac{\eu^{3x}\,\eu^{-x+1}}{\eu^{x}}$ et
$\ln\!\left(\dfrac{\eu^2\sqrt{\eu}}{\eu^{-3}}\right)$.
\end{exercise}

\begin{exercise}[$\star$]\label{exo:g12:exp:2}
Résoudre dans $\R$~:
\[
\text{(a) } \eu^{2x} - 3\eu^{x} + 2 = 0;
\qquad
\text{(b) } \ln(x - 1) + \ln(x + 2) = \ln 4 .
\]
\end{exercise}

\begin{exercise}[$\star$]\label{exo:g12:exp:3}
Calculer les limites~:
\[
\lim_{x\to+\infty} \frac{\eu^x - x^2}{\eu^x + 1}, \qquad
\lim_{x\to+\infty} \frac{\ln x}{\sqrt x}, \qquad
\lim_{x\to0^+} x^2 \ln x, \qquad
\lim_{x\to+\infty} \bigl(x - \ln x\bigr).
\]
\end{exercise}

\begin{exercise}[$\star\star$]\label{exo:g12:exp:4}
Étudier la \omterm{def:g11:func:function}{fonction} $f(x) = x\,\eu^{-x}$ sur $\R$~: variations, limites,
\omterm{def:g12:deriv:extremum}{extremum}~; montrer que sa courbe a un \omterm{def:g12:deriv:inflection}{point d'inflexion} et donner ses
\omterm{def:g10:coordgeom:system}{coordonnées}.
\end{exercise}

\begin{exercise}[$\star\star$]\label{exo:g12:exp:5}
Montrer que pour tout $x > -1$, $\ln(1 + x) \leq x$, avec égalité seulement
en $x = 0$. En déduire que pour tout $n \geq 1$,
\[
\left(1 + \frac{1}{n}\right)^n \leq \eu \leq \left(1 - \frac{1}{n+1}\right)^{-(n+1)} .
\]
\end{exercise}

\begin{exercise}[$\star\star$]\label{exo:g12:exp:6}
Un capital $C_0$ est placé au taux annuel de $3\,\%$, intérêts composés
chaque année.
\begin{enumerate}
  \item Exprimer le capital $C_n$ après $n$ années.
  \item Au bout de combien d'années le capital double-t-il~? Donner la
        réponse exacte avec $\ln$, puis une valeur numérique.
  \item Comparer avec l'\omterm{def:g10:numbers:approx}{approximation} «~$70$ divisé par le taux en
        pourcent~» utilisée par les banquiers.
\end{enumerate}
\end{exercise}

\begin{exercise}[$\star\star$]\label{exo:g12:exp:7}
Résoudre l'inéquation $\eu^{2x} - \eu^{x+1} > 0$, puis l'inéquation
$\ln(x^2 - 1) \leq \ln(x + 5)$.
\end{exercise}

\begin{exercise}[$\star\star\star$]\label{exo:g12:exp:8}
Soit $f(x) = \dfrac{\eu^x}{x}$ pour $x > 0$.
\begin{enumerate}
  \item Étudier les variations de $f$ sur $\intoo{0}{+\infty}$ et donner
        son \omterm{def:g10:functions:extrema}{minimum}.
  \item Pour quelles valeurs de $k$ l'\omterm{def:g10:algebra:equation}{équation} $\eu^x = kx$ a-t-elle $0$,
        $1$ ou $2$ solutions dans $\intoo{0}{+\infty}$~?
\end{enumerate}
\end{exercise}

\begin{exercise}[$\star\star\star$]\label{exo:g12:exp:9}
Pour $n \geq 1$, soit $u_n = \left(1 + \frac1n\right)^n$.
\begin{enumerate}
  \item En utilisant l'\cref{exo:g12:exp:5}, montrer que $(u_n)$ est
        \omterm{def:g12:seq:bounded}{majorée} par~$\eu$.
  \item Montrer que $\ln u_n = n \ln\left(1 + \frac1n\right) \to 1$, et en
        déduire que $u_n \to \eu$.
        (Indication~: reconnaître un taux d'accroissement de $\ln$ en $1$.)
\end{enumerate}
\end{exercise}

\section{Problème~: le logarithme dompte le monde}

\begin{problem}[{Devoir du week-end --- temps de doublement et règle
de 72, les échelles des séismes, des acides et des pianos, et la
constante $\eu$ qui laisse partout ses empreintes}]
\label{pb:g12:exp:1}
Tout ce qui croît d'un \emph{pourcentage} fixe croît
exponentiellement --- l'épargne, les bactéries, les épidémies ---
et tout ce qui s'étend sur trop de puissances de dix pour être saisi
--- énergies des séismes, acidités, intensités sonores --- se
laisse dompter par un \omterm{def:g12:exp:ln}{logarithme}. Ce problème calcule des temps de
doublement et démasque la règle de 72 des banquiers, lit les
échelles logarithmiques du monde, et recueille les empreintes que le
nombre $\eu$ laisse sur toutes les scènes de crime
(\cref{prop:g12:exp:funceq}, \cref{thm:g12:exp:growth},
\cref{met:g12:exp:equations}).

\medskip
\noindent\textbf{Partie I --- Aisance.}
\begin{enumerate}
  \item Résoudre~: $\eu^{2x} = 5$~; \quad $\ln(3x - 1) = 2$~; \quad
        $\eu^{2x} - 3\eu^x + 2 = 0$ (un trinôme déguisé).
  \item Simplifier~: $\dfrac{\ln 8}{\ln 2}$~; \quad
        $\ln\!\left(\eu^3 \sqrt{\eu}\right)$~; \quad
        $\eu^{\ln 5 - \ln 2}$.
  \item Démontrer l'inégalité mère~: $\eu^x \geq 1 + x$ pour tout
        réel $x$ (étudier $f(x) = \eu^x - 1 - x$), et en déduire
        $\ln(1 + u) \leq u$ pour $u > -1$.
  \item Calculer $\lim_{x \to +\infty} x^2 \eu^{-x}$,
        $\lim_{x \to +\infty} \frac{\ln x}{x}$
        (\cref{thm:g12:exp:growth}), et
        $\lim_{x \to 0} \frac{\eu^x - 1}{x}$ (y reconnaître une
        \omterm{def:g12:deriv:derivative}{dérivée}).
  \item Étudier $f(x) = x \ln x$ sur $\intoo{0}{+\infty}$~:
        variations, \omterm{def:g10:functions:extrema}{minimum}, et limite en $0^+$ (on admet que
        $x \ln x \to 0$). Cette petite \omterm{def:g11:func:function}{fonction} mesure l'information
        et l'entropie tout au long des volumes universitaires.
\end{enumerate}

\medskip
\noindent\textbf{Partie II --- Temps de doublement et règle de
72.}
\begin{enumerate}[resume]
  \item Une épargne croît de $3\,\%$ par an. Résoudre
        $1.03^n = 2$~: en combien de temps le capital
        double-t-il~?
  \item Les banquiers estiment le temps de doublement par
        $\frac{72}{\text{taux en }\%}$. Tester la règle à $3\,\%$,
        $6\,\%$ et $9\,\%$ face à la valeur exacte
        $\frac{\ln 2}{\ln(1 + r)}$, puis l'expliquer~: pour $r$
        petit, $\ln(1 + r) \approx r$ (l'inégalité de la question 3
        en donne la moitié), si bien que la constante exacte est
        $100 \ln 2 \approx 69.3$ --- pourquoi les banquiers
        préfèrent-ils $72$~?
  \item Une bactérie se divise toutes les $20$ minutes~:
        $p(t) = 2^{t/20}$ ($t$ en minutes). La réécrire sous la
        forme $\eu^{\lambda t}$, puis calculer $p$ au bout de $24$
        heures. Que prouve la réponse (plus de $10^{21}$) sur le
        modèle --- et qu'est-ce qui arrête les colonies réelles~?
  \item La caféine quitte l'organisme avec une demi-vie d'environ
        $5$ heures. Après un café de $100$ mg pris à 15 h, combien
        en reste-t-il à 23 h~? À quelle heure passe-t-on sous
        $10$ mg~? (L'insomnie a un \omterm{def:g12:exp:ln}{logarithme}.)
  \item Un euro à $100\,\%$ d'intérêt annuel~: calculer le capital
        en fin d'année pour une capitalisation annuelle, mensuelle
        puis quotidienne, et donner le plafond que
        l'\cref{exo:g12:exp:9} a démontré infranchissable. Quelle
        constante célèbre est la limite de la pure cupidité~?
  \item Pourquoi les scientifiques tracent-ils
        $\ln(\text{nombre de cas})$ en \omterm{def:g11:func:function}{fonction} du temps lors de la
        phase initiale d'une épidémie~? Que révèle une droite sur ce
        \omterm{def:g10:functions:graph}{graphique}, et que mesure son \omterm{def:g10:reffunc:affine}{coefficient directeur}~?
\end{enumerate}

\medskip
\noindent\textbf{Partie III --- Les échelles logarithmiques du
monde.}
(On écrit $\log_{10} x = \frac{\ln x}{\ln 10}$.)
\begin{enumerate}[resume]
  \item Niveau sonore en décibels~: $L = 10 \log_{10}(I/I_0)$. Une
        conversation mesure $60$ dB, un concert de rock $120$ dB~:
        par quel facteur les intensités sonores diffèrent-elles~?
  \item Les magnitudes des séismes augmentent de $1$ lorsque
        l'amplitude sismique est multipliée par $10$, et l'énergie
        libérée varie comme amplitude$^{3/2}$. Comparer un séisme de
        magnitude $5$ et un de magnitude $7$~: rapport des
        amplitudes, puis rapport des énergies.
  \item Chimie~: $\text{pH} = -\log_{10}[\mathrm{H^+}]$. Le jus de
        citron a un pH de $2$, le lait un pH de $6.5$~: quel est le
        rapport de leurs concentrations en acide~?
  \item Musique~: chaque octave double la \omterm{def:g10:stats:series}{fréquence}. Du la le plus
        grave du piano ($27.5$ Hz) à son do le plus aigu
        ($4\,186$ Hz), combien d'octaves le clavier couvre-t-il~? Et
        pourquoi nos sens --- ouïe, vue, perception des secousses
        --- préfèrent-ils les échelles logarithmiques~? (Une
        phrase.)
  \item La règle à calcul, calculatrice des ingénieurs pendant
        $350$ ans~: deux réglettes graduées de sorte que la
        \emph{longueur} jusqu'au trait $x$ soit proportionnelle à
        $\ln x$. Expliquer comment le glissement d'une réglette le
        long de l'autre multiplie les nombres, et nommer l'identité
        de la \cref{prop:g12:exp:lnalg} qui fait le travail.
\end{enumerate}

\medskip
\noindent\textbf{Partie IV --- Les empreintes de $\eu$.}
\begin{enumerate}[resume]
  \item D'après l'\cref{exo:g12:exp:8}~: l'\omterm{def:g10:algebra:equation}{équation} $\eu^x = kx$
        ($x > 0$) a $0$, $1$ ou $2$ solutions selon la position de
        $k$ par rapport à un seuil. Réénoncer le résultat --- et
        vérifier le fait géométrique qui le sous-tend~: la droite
        $y = \eu x$ est exactement la \omterm{def:g12:deriv:tangent}{tangente} à l'exponentielle
        passant par l'origine.
  \item La constante du presque-raté~: avec $n$ billets de loterie
        gagnants avec la \omterm{def:g10:proba:distribution}{probabilité} $\frac1n$ chacun,
        $\P(\text{aucun gain}) = \left(1 - \frac1n\right)^n$.
        Calculer sa limite en passant par
        $n \ln\!\left(1 - \frac1n\right)$ (utiliser la limite
        usuelle $\frac{\ln(1 + u)}{u} \to 1$), et faire le lien avec
        le mystérieux $0.368$ du \cref{pb:g11:binom:1}.
  \item La règle des $37\,\%$~: pour choisir le meilleur parmi $n$
        candidats reçus dans un ordre aléatoire (sans retour en
        arrière), la stratégie optimale consiste à écarter les
        $\frac{n}{\eu} \approx 37\,\%$ premiers, puis à retenir le
        premier candidat meilleur que tous les précédents --- ce qui
        réussit avec une \omterm{def:g10:proba:distribution}{probabilité} voisine de $\frac1\eu$. D'où
        viennent les deux $\frac1\eu$ des questions 18 et 19 ---
        énoncer le mécanisme commun (de nombreux \omterm{def:g11:prob:model}{événements}
        indépendants, chacun peu probable) --- et calculer
        $\frac1\eu$ à trois décimales.
  \item Pour finir --- le portrait de $\eu$~: le plafond de la
        capitalisation (question 10)~; la base dont la \omterm{def:g12:deriv:tangent}{tangente} en
        $0$ a un \omterm{def:g10:reffunc:affine}{coefficient directeur} exactement égal à $1$
        (questions 3 et 4)~; la croissance qu'aucun polynôme ne
        rattrape (question 4)~; la constante des presque-ratés et de
        l'arrêt optimal (questions 18 et 19)~; et son inverse,
        $\ln$, qui change les produits en sommes (question 16) et
        les décennies de puissances en centimètres (partie III). Une
        phrase pour chacun.
\end{enumerate}
\end{problem}
